\documentclass[10pt]{article}
\usepackage[a4paper, top=0.5in, bottom=0.5in, left=0.5in, right=0.5in,
            includeheadfoot]{geometry}
\usepackage[T1]{fontenc}
\usepackage[utf8]{inputenc}
\usepackage{setspace}
\usepackage{parskip}
\usepackage{xurl}
\usepackage{hyperref}
\usepackage{qrcode}
\usepackage{fancyhdr}
\usepackage{lastpage}
\hypersetup{breaklinks=true, colorlinks=true, linkcolor=black, citecolor=black, urlcolor=blue}

\setstretch{1.15}
\setlength{\parskip}{4pt}

% ── Header / Footer ────────────────────────────────────────────────────────────
\pagestyle{fancy}
\fancyhf{}
\renewcommand{\headrulewidth}{0.4pt}
\renewcommand{\footrulewidth}{0.4pt}
\fancyhead[L]{\small\textit{CH4217 --- Research Methodology}}
\fancyhead[C]{\small\textit{Birbal Sahni and the Architecture of Scientific Sovereignty}}
\fancyhead[R]{\small\textit{Shuvam Banerji Seal}}
\fancyfoot[L]{\small IISER Kolkata $\cdot$ Department of Chemical Sciences}
\fancyfoot[C]{\small Page \thepage\ of \pageref{LastPage}}
\fancyfoot[R]{\small April 12, 2026}

\fancypagestyle{titlepage}{%
  \fancyhf{}
  \renewcommand{\headrulewidth}{0pt}
  \renewcommand{\footrulewidth}{0.4pt}
  \fancyfoot[L]{\small IISER Kolkata $\cdot$ Department of Chemical Sciences}
  \fancyfoot[C]{\small Page \thepage\ of \pageref{LastPage}}
  \fancyfoot[R]{\small April 12, 2026}
}

\title{\textbf{Birbal Sahni and the Architecture of Scientific Sovereignty in India}}
\author{%
  Shuvam Banerji Seal\\[2pt]
  \small Roll No.: 22MS076\quad 4th Year BS-MS\\
  \small Department of Chemical Sciences, IISER Kolkata\\
  \small Course: CH4217 ``Research Methodology''
}
\date{April 12, 2026}

\begin{document}
\maketitle
\thispagestyle{titlepage}

\begin{center}
\small\textit{Scan for digital companion resources:}\\[4pt]
\begin{tabular}{c@{\hspace{1.2cm}}c}
\qrcode[height=2.2cm]{https://github.com/Shuvam-Banerji-Seal/CH4217_Research_Methodology} &
\qrcode[height=2.2cm]{https://shuvam-banerji-seal.github.io/CH4217_Research_Methodology/} \\[2pt]
\small Repository & \small Live Website
\end{tabular}
\end{center}

\vspace{2pt}\hrule\vspace{6pt}


\section*{Introduction}

Birbal Sahni's career poses a question that exceeds conventional biography: how does a scientist convert rigorous personal scholarship into durable public capacity, and how does that conversion reshape a discipline whose institutional foundations have yet to be built? In late colonial India,\cite{science_state_nation} natural history frequently served the administrative and extractive demands of imperial survey, whereas disciplinary interpretation of fossil material remained anchored in metropolitan authority; however, Sahni methodically altered that structure by fusing paleobotanical analysis, evolutionary reasoning, and institution-building strategy into a single coherent research program.\cite{wiki_birbal,palaeobot_bio} His career consequently marks a transition of considerable analytical significance: from data collection under empire to autonomous knowledge production within an emerging scientific nation, and this transition carries implications that remain relevant for research methodology today.

This essay develops that argument through a sequential analysis of his biographical origins, formal education, central research contributions, most notable achievements, and enduring social impact. It begins with the geographical and familial foundations that cultivated his scientific temperament,\cite{from_lahore_lucknow} then traces how his methodological innovations in paleobotany conferred wider epistemological legitimacy on an entire field, and finally evaluates why his institutional legacy constitutes an instructive operational model for research training in the contemporary Indian context.\cite{bsip_history}


\section*{Origin, Birth, and Family Background}

Birbal Sahni was born on 14 November 1891 in Bhera, a town in Western Punjab then in British India and now in Pakistan.\cite{wiki_birbal} The geographical significance of this birthplace exceeds mere biographical detail; Bhera lies in close proximity to the Salt Range, a geological formation so densely fossiliferous that it was sometimes described as a ``Museum of Geology,'' and this landscape provided the adolescent Sahni with an unmediated encounter with stratified deep-time records before formal academic specialization had begun.\cite{from_lahore_lucknow} The natural archive surrounding his childhood operated as an early pedagogical environment of unusual quality; it cultivated a geological curiosity that would later anchor his most productive scientific work.

The household in which Sahni was raised further amplified this formative setting. His father, Lala Ruchi Ram Sahni, held a professorship in chemistry at Government College Lahore\cite{ruchi_ram} and participated actively in a reformist scientific culture shaped by rational pedagogy and civic public education; Ruchi Ram delivered popular science lectures to rural audiences using ``Magic Lantern'' projection technology, thereby encoding in the family ethos a conviction that scientific knowledge carried social obligations.\cite{punjab_science} Accounts of the household describe an environment where scientific discussion, field observation, and political consciousness coexisted without contradiction, and this family-learning ecology appears to have cultivated Birbal Sahni's subsequent commitment to both rigorous empirical inquiry and knowledge democratization.\cite{from_lahore_lucknow} The family's association with Brahmo Samaj values, which emphasized rational critique over scriptural authority and social reform over cultural passivity, further reinforced an anti-obscurantist intellectual orientation that remained visible throughout Sahni's professional conduct; consequently, when he later challenged received European classifications with evidence drawn from Indian material, he was extending a disposition inherited from a reformist domestic context, not simply asserting disciplinary novelty.


\section*{Education and Intellectual Formation}

Sahni received foundational scientific training at Government College Lahore, where he completed botanical and geological coursework before proceeding to Cambridge.\cite{from_lahore_lucknow} His admission to Emmanuel College around 1914 and his placement under Professor Albert Charles Seward, then the leading paleobotanist in the English-speaking world, constituted the decisive inflection point in his methodological development. Under Seward's guidance, Sahni absorbed comparative morphology, stratigraphic reasoning, and taxonomic discipline at a depth that allowed him to evaluate fossil evidence as an interconnected biological and geological system rather than as a catalogue of isolated specimens; Seward later declined to study fossil collections forwarded by the Geological Survey of India, insisting that this material should be interpreted by his ``young pupil'' who held prior interpretive rights over his own nation's heritage.\cite{from_lahore_lucknow}

The most consequential dimension of Sahni's training, however, emerged not in Cambridge itself but in the strategic posture he adopted upon his return to India in 1919.\cite{wiki_birbal} Many colonial-era scholars remained embedded in metropolitan validation structures even after returning home, reproducing intellectual dependence through habitual deference to European journals and institutions; Sahni pursued a distinct return-with-purpose strategy, applying global methodological standards rigorously to Indian material while simultaneously asserting interpretive priority for Indian scientists.\cite{science_state_nation} This repositioning transformed education into a sovereignty instrument: the acquisition of international methodological competence became not a credential to be displayed for foreign audiences but a tool for building domestic intellectual infrastructure, and this model of purposive international training followed by deliberate institutional reinvestment remains instructive for contemporary research policy.\cite{from_lahore_lucknow}


\section*{Research Contributions: Method, Evidence, and Synthesis}

Sahni's principal research contribution lies in transforming Indian paleobotany from descriptive fossil cataloguing into evidence-integrated biological interpretation. Before this methodological shift, plant fossils in India were routinely treated as stratigraphic markers within geological survey routines, their biological significance systematically underweighted; Sahni reoriented the field by linking morphological analysis, depositional context, and evolutionary inference into a unified interpretive framework, and this integration expanded both the explanatory scope and the disciplinary credibility of paleobotanical research within the broader earth sciences.\cite{palaeobot_bio}

His work on the Gondwana flora, particularly \textit{Glossopteris}, most clearly illustrates the structure of that transformation. By analyzing fossil distributions in relation to glacial and periglacial depositional contexts, Sahni challenged simplified floral succession assumptions and demonstrated that \textit{Glossopteris} was adapted to cold-temperate and even periglacial conditions rather than appearing only after the retreat of Paleozoic ice sheets;\cite{regionalism_mobilism} the implication was systemic, repositioning Indian plant fossils as central evidence in continental-drift debates rather than as local curiosities in museum catalogs. His 1938 meeting at the Indian Science Congress with South African geologist Alexander du Toit, the preeminent advocate of the Gondwanaland concept,\cite{from_lahore_lucknow} consolidated this biogeographic argument through direct empirical cross-referencing: the near-identity of Gondwana floras across India, South Africa, South America, and Australia provided biological proof of shared geological ancestry that hypothetical land bridges could not adequately explain.\cite{regionalism_mobilism}

A second major contribution was whole-plant reconstruction from anatomically disarticulated organs, performed most productively in the silicified fossil beds of the Rajmahal Hills of Bihar.\cite{palaeobot_bio} Before this work, different structural components of the same extinct plant were routinely described as taxonomically distinct genera because field collection produced them in isolation; Sahni's meticulous anatomical correlations demonstrated that the stem \textit{Bucklandia}, the leaf \textit{Ptilophyllum}, and the reproductive structure \textit{Williamsonia} constituted a single organism, subsequently reconstructed as \textit{Williamsonia sewardiana}.\cite{from_lahore_lucknow} This organ-correlation method reduced taxonomic fragmentation and shifted interpretive emphasis toward organism-level reconstruction, moving the field toward developmental coherence and functional biological reasoning.\cite{palaeobot_bio}

His late work, culminating in the 1948 paper introducing the \textit{Pentoxyleae} (now \textit{Pentoxylales}), further demonstrates the epistemic discipline that characterized his approach throughout: the group displayed a trait combination, including polystelic vascular anatomy reminiscent of Paleozoic \textit{Medullosae}, conifer-like secondary wood, and fleshy cone-like reproductive organs without close analogue in any known group, that resisted neat linear taxonomic placement.\cite{pentoxylales_affinities} Sahni treated this morphological anomaly as evidence to be explained rather than noise to be suppressed; in methodological terms, this stance reflects high epistemic discipline, because it required theory to accommodate inconvenient data rather than permitting selective data pruning to preserve theoretical elegance.\cite{from_lahore_lucknow}


\section*{Notable Achievements and Institutional Leadership}

Sahni's notable achievements operated across two registers that reinforced one another: personal scientific recognition and deliberate institutional construction. His election as a Fellow of the Royal Society in 1936 marked him as the first Indian botanist to receive that distinction,\cite{wiki_birbal} signalling international acknowledgment of his authority to interpret India's paleobotanical record on equal terms with metropolitan scientists; yet the more consequential outcome was how he directed that authority, consistently investing it in building domestic research infrastructure rather than exploiting it for individual prestige accumulation.

In May 1946, Sahni convened eight researchers in Lucknow to sign the Memorandum of Association establishing the Palaeobotanical Society; the Institute of Palaeobotany was formally inaugurated on 10 September 1946, with Sahni as its founding honorary director, on the basis of his own and Savitri Sahni's personal endowment of their library, fossil collections, and an initial financial trust.\cite{bsip_history} This institutional act aligned with the broader transition from colonial extraction science to development-oriented national science in the years immediately surrounding Independence, and the institute functioned simultaneously as a research center, an archival repository, and a sovereignty institution, because it anchored interpretive labor, disciplinary training, and fossil curation physically within India.\cite{bsip_wiki}

Prime Minister Jawaharlal Nehru attended the laying of the institute's foundation stone on 3 April 1949, regarding the institution as a cornerstone of the new nation's intellectual architecture.\cite{bsip_history} Sahni suffered a fatal heart attack seven days later, on 10 April 1949, at the age of 57.\cite{wiki_birbal} The continuity of his vision depended substantially on Savitri Sahni's administrative and scientific stewardship; she oversaw completion of the main building by late 1952 and secured internationally respected scientists, including Professor T.~M.~Harris of Reading University and Professor O.~A.~H\o eg of Oslo, as advisory directors during the institution's formative transition period.\cite{from_lahore_lucknow} That a fledgling institute survived founder loss without significant loss of research direction confirms that Sahni's contribution cannot be reduced to charismatic personal leadership; he had designed a transferable institutional architecture capable of sustaining itself across personnel changes.


\section*{Social Impact and Public Significance}

The social impact of Sahni's work operated across scientific, educational, and civic registers simultaneously. Scientifically, he established a discipline that trained successive generations of Indian researchers to engage with domestic fossil archives at internationally competitive methodological standards, thereby reversing the directional dependency in which Indian material was routinely sent abroad for interpretation.\cite{from_lahore_lucknow} Educationally, he advanced a research culture in which field evidence, analytical rigor, and conceptual synthesis were taught as an integrated craft rather than as discrete procedural skills, and his tenure as Professor and Head of the Botany Department at Lucknow University from 1921 until his death produced a cohort of researchers who extended this culture well beyond his own research agenda.\cite{bsip_history}

His polymathic engagements widened the public relevance of this scientific culture further still. His 1945 monograph \textit{The Technique of Casting Coins in Ancient India} reconstructed the precision casting methods of Yaudheya tribal coinage at Khokhra-Kot through direct morphological analysis of coin moulds,\cite{palaeobot_bio} earning him the Nelson Wright Medal from the Numismatic Society of India; separately, he used plant anatomical evidence from charred wood at Harappan sites to infer sophisticated Indus Valley trade networks with Himalayan timber suppliers,\cite{wiki_birbal} demonstrating that rigorous evidence methods could illuminate economic and social history as effectively as biological evolution. These interdisciplinary contributions produced a broader social legitimacy for evidence-based science, because they showed that analytical standards developed in paleobotany could migrate across disciplinary boundaries without loss of precision or interpretive power.

His international collaborations further positioned Indian paleobotany within transregional evidence networks. His decade-long mentoring of Chinese scientist Xu Ren (Hsu Jen), who arrived in Lucknow in 1944 carrying fossil specimens airlifted over the eastern Himalayas during the Second World War,\cite{trans_himalayan_article} concluded when Xu Ren earned what was likely the first doctoral degree awarded to a Chinese scientist by an Indian university in 1946\cite{trans_himalayan_pdf} and subsequently established the Institute of Botany at the Chinese Academy of Sciences in 1959 on the institutional model Sahni had developed. This demonstrated that the Lucknow model was exportable and that Indian science could function as a generative network node rather than a perpetual recipient periphery within global knowledge structures.\cite{trans_himalayan_pdf}


\section*{Contemporary Relevance of the Sahni Legacy}

The contemporary Birbal Sahni Institute of Palaeosciences extends his framework into methodological domains that were unavailable to him: palynology, stable isotope analysis, biomarker geochemistry, dendrochronology, and ancient DNA research.\cite{dna_lab,palaeoscience_today} Under senior scientist Niraj Rai, the Institute's ancient DNA laboratory now applies genomic analysis to human remains from Harappan sites, Mesolithic Ganga plains populations, and South Indian megalithic cultures, reconstructing the demographic and population-genetic history of the subcontinent through direct molecular evidence rather than inference from modern proxy distributions.\cite{dna_lab} These expansions represent disciplined methodological continuity rather than institutional drift: fragmentary traces are still assembled through layered evidence logic to reconstruct historical processes, and the epistemological commitment to evidence-over-assertion that Sahni embodied remains the governing standard across all research programs.\cite{sponsored_projects,isti_bsip}

This continuity carries a pointed methodological lesson for research training in India today. Durable scholarship requires more than publication output; it requires archive stewardship, intergenerational mentoring, instrument investment, and institution-level memory systems that preserve methodological standards across the inevitable attrition of individual researchers.\cite{bsip_history,science_state_nation} Sahni's career, read in this light, offers an operational model of how scientific excellence can be socially embedded and historically sustained rather than remaining contingent on individual genius: the institute he built has outlasted him by more than seven decades, and its continued international productivity affirms the strategic wisdom of deliberate institution-building investment.


\section*{Conclusion}

Birbal Sahni's biography begins in Bhera with a precise birth record and a scientifically active family context, but its larger significance lies in what he built from those beginnings: through advanced education, rigorous research contribution, notable disciplinary achievements, and long-range institution design, he helped convert paleobotany in India from a peripheral descriptive practice into a sovereign knowledge field with global relevance; his social impact persists because he treated method, training, and institution as interdependent components of scientific life rather than as separable career choices.\cite{wiki_birbal,from_lahore_lucknow,bsip_history}

For a course in research methodology, Sahni's legacy demonstrates a demanding but fruitful principle: evidence gains historical force when analytical precision is matched by organizational foresight. Scientific nation-building, in this perspective, is neither rhetorical aspiration nor policy slogan; it is the cumulative result of disciplined inference, strategic collaboration, and institutions designed to outlast their founders.\cite{science_state_nation,trans_himalayan_article,isti_bsip}

Yet one must ask whether the epistemic sovereignty Sahni laboured to construct is being sustained or systematically hollowed. India's gross expenditure on research and development remains near 0.64\% of GDP, a figure that has stagnated for over a decade while peer economies such as South Korea (4.9\%), Israel (5.4\%), and China (2.4\%) accelerate their investment trajectories, and the resultant resource deficit propagates a compounding feedback loop of institutional underfunding, laboratory obsolescence, and investigator attrition.\cite{dst_rd_stats,nature_brain_drain} Between 2018 and 2024 the annual outflow of Indian science and engineering doctoral candidates to OECD destination countries rose by approximately 55\%, a human-capital haemorrhage that constitutes a measurable erosion of precisely the domestic knowledge-reproduction capacity Sahni considered foundational to scientific sovereignty.\cite{nsf_sed,oecd_migration} Meanwhile, curriculum revision committees have introduced astrology modules and mythological origin claims into university science syllabi,\cite{ugc_astrology,curr_pseudo} a substitution of pseudoscientific assertion for evidence-based reasoning that corrodes precisely the anti-obscurantist temperament Sahni inherited from Ruchi Ram's rationalist pedagogy and embodied throughout his professional career. When funding stagnates, when trained researchers emigrate because domestic opportunity structures contract, and when policy frameworks conflate civilizational pride with empirical method, the institutional architecture of scientific sovereignty does not merely weaken; it undergoes a phase transition toward irreversibility. So one must confront the question directly: does the India of today continue the legacy Birbal Sahni built, stone by disciplined stone, or does it dismantle that legacy while claiming to honour it?

\begin{thebibliography}{99}

\bibitem{wiki_birbal}
Wikipedia contributors, ``Birbal Sahni,''
\url{https://en.wikipedia.org/wiki/Birbal_Sahni} (accessed April 12, 2026).

\bibitem{palaeobot_bio}
International Organisation of Palaeobotany, ``Sahni, Birbal (1891--1949),''
\url{https://www.palaeobotany.org/index.php/palaeobotanist-biographies/sahni-birbal-1891-1949/} (accessed April 12, 2026).

\bibitem{science_state_nation}
Deepak Kumar, ``Science, State and Nation,'' in \textit{Science and the Raj},
Cambridge University Press,
\url{https://resolve.cambridge.org/core/services/aop-cambridge-core/content/view/FA142C30EB3A60B8CB768216DD0EC947/9781139053426c6_p169-210_CBO.pdf/science_state_and_nation.pdf}
(accessed April 12, 2026).

\bibitem{from_lahore_lucknow}
``From Lahore to Lucknow: The Foundational Legacy of Birbal Sahni in Indian
Paleobotany,'' source manuscript provided in course workspace (accessed April 12, 2026).

\bibitem{bsip_history}
Birbal Sahni Institute of Palaeosciences, ``Institute History,''
\url{https://bsip.res.in/en/bsip_institute_history} (accessed April 12, 2026).

\bibitem{punjab_science}
``Birbal Sahni and His Father Ruchi Ram: Science in Punjab Emerging from the
Shadows of the Raj,''
\url{https://www.academia.edu/92746487/Birbal_Sahni_and_His_Father_Ruchi_Ram_Science_in_Punjab_Emerging_from_the_Shadows_of_the_Raj}
(accessed April 12, 2026).

\bibitem{ruchi_ram}
Wikipedia contributors, ``Ruchi Ram Sahni,''
\url{https://en.wikipedia.org/wiki/Ruchi_Ram_Sahni} (accessed April 12, 2026).

\bibitem{regionalism_mobilism}
N.~Oreskes (ed.), ``Regionalism and the Reception of Mobilism: South Africa,
India, and South America from the 1920s through the Early 1950s,'' in
\textit{The Continental Drift Controversy}, Cambridge University Press,
\url{https://www.cambridge.org/core/books/continental-drift-controversy/regionalism-and-the-reception-of-mobilism-south-africa-india-and-south-america-from-the-1920s-through-the-early-1950s/7C14C2C0F8BA04DF8829B90B9AF06361}
(accessed April 12, 2026).

\bibitem{pentoxylales_affinities}
``The Affinities of the Pentoxyleae,'' \textit{Palaeobotanist} archive,
\url{http://14.139.63.228:8080/pbrep/bitstream/123456789/1081/1/Pb2829_207.pdf}
(accessed April 12, 2026).

\bibitem{trans_himalayan_pdf}
A.~Ghosh, ``Trans-Himalayan Science in Mid-Twentieth Century China and India:
Birbal Sahni, Hsu Jen, and a Pan-Asian Paleobotany,''
\url{https://www.cambridge.org/core/services/aop-cambridge-core/content/view/59215D119B63EE5F1B00BF4DBB05F51E/S1479591421000292a.pdf/transhimalayan_science_in_midtwentieth_century_china_and_india_birbal_sahni_hsu_jen_and_a_panasian_paleobotany.pdf}
(accessed April 12, 2026).

\bibitem{trans_himalayan_article}
A.~Ghosh, ``Trans-Himalayan Science in Mid-Twentieth Century China and India,''
\textit{International Journal of Asian Studies}, Cambridge Core,
\url{https://www.cambridge.org/core/journals/international-journal-of-asian-studies/article/transhimalayan-science-in-midtwentieth-century-china-and-india-birbal-sahni-hsu-jen-and-a-panasian-paleobotany/59215D119B63EE5F1B00BF4DBB05F51E}
(accessed April 12, 2026).

\bibitem{bsip_wiki}
Wikipedia contributors, ``Birbal Sahni Institute of Palaeosciences,''
\url{https://en.wikipedia.org/wiki/Birbal_Sahni_Institute_of_Palaeosciences}
(accessed April 12, 2026).

\bibitem{isti_bsip}
India Science, Technology and Innovation Portal, ``Birbal Sahni Institute of
Palaeosciences (BSIP), Lucknow,''
\url{https://www.indiascienceandtechnology.gov.in/organisations/ministry-and-departments/department-science-technology-dst/birbal-sahni-institute}
(accessed April 12, 2026).

\bibitem{palaeoscience_today}
Birbal Sahni Institute of Palaeosciences, \textit{Palaeoscience Today},
Oct.--Dec.\ 2025 issue,
\url{https://bsip.res.in/assets-new/assets/paleoscienceToday/Palaeoscience%20Today,%20Oct.%20-%20Dec.%202025.pdf}
(accessed April 12, 2026).

\bibitem{sponsored_projects}
Birbal Sahni Institute of Palaeosciences, ``Sponsored Projects,''
\url{https://www.bsip.res.in/en/bsip_sponsored_project} (accessed April 12, 2026).

\bibitem{dna_lab}
Birbal Sahni Institute of Palaeosciences, ``Ancient DNA Laboratory,''
\url{https://www.bsip.res.in/en/bsip_dna_lab} (accessed April 12, 2026).

\bibitem{dst_rd_stats}
Press Information Bureau, Government of India,
\textit{Research and Development Statistics at a Glance 2022--23},
\url{https://pib.gov.in/PressReleasePage.aspx?PRID=1951750}
(accessed April 12, 2026).

\bibitem{nature_brain_drain}
G.~Trivedi, ``India's Brain Drain: Why Scientists Are Leaving and What It
Means,'' \textit{Nature India}, 2024,
\url{https://doi.org/10.1038/d44151-024-00050-4} (accessed April 12, 2026).

\bibitem{nsf_sed}
National Science Foundation, \textit{Survey of Earned Doctorates 2023},
Table~24: Non-U.S.\ citizen doctorate recipients by country of origin,
\url{https://ncses.nsf.gov/surveys/earned-doctorates/2023}
(accessed April 12, 2026).

\bibitem{oecd_migration}
OECD, \textit{International Migration Outlook 2024}, Chapter~3: Highly skilled
migration flows, \url{https://doi.org/10.1787/29f23e9d-en} (accessed April 12, 2026).

\bibitem{ugc_astrology}
University Grants Commission, notification on ``Jyotir Vigyan'' (Vedic
Astrology) as a university discipline; coverage and commentary in
\textit{The Wire},
\url{https://thewire.in/education/ugc-astrology-vedic-science-curriculum-universities}
(accessed April 12, 2026).

\bibitem{curr_pseudo}
\textit{The Wire}, ``UGC's Push for Vedic Science and Astrology in University
Curricula,''
\url{https://thewire.in/education/ugc-astrology-vedic-science-curriculum-universities}
(accessed April 12, 2026).

\end{thebibliography}

\end{document}